\documentclass[12pt]{article}
\usepackage[utf8]{inputenc}
\usepackage[T1]{fontenc}
\usepackage{hyperref}
\title{Bidayuh Serian Language Resources for K-12 Curriculum Development}
\author{Compiled by Claude Sonnet 5.5 (claude-sonnet-5-5) for Andr\'as Kornai}
\date{}
\begin{document}
\maketitle

\section{Introduction}
Bidayuh Serian (ISO 639-3 sdo; also called Bukar-Sadong, Buka, Sadong or Serian Bidayuh) is an Austronesian (Land Dayak) language of Sarawak, Malaysia, and West Kalimantan, Indonesia, with about 49,000 speakers. Searches found only general descriptions, a few web resources and news about school plans. No textbook or primer specific to this variety turned up.

\section{Primary Educational Resources}
No primers or textbooks for Bukar-Sadong were found. News reports say a Bidayuh subject for Sarawak schools has been proposed, but the reports do not say which dialect it would use.
\begin{itemize}
\item Borneo Post (2025): the state minister said Bidayuh may be introduced as a school subject, as an optional subject under the Ministry of Education's ethnic language curriculum. \url{https://www.theborneopost.com/2025/08/23/sagah-bidayuh-language-may-be-introduced-as-school-subject-next-year/}
\item Serian Division news page carrying the same story. \url{https://serian.sarawak.gov.my/web/subpage/news_view/1075}
\item UNIMAS record on Bidayuh heritage-language playschools and kindergartens in Sarawak. \url{https://ir.unimas.my/id/eprint/19669/}
\item UNIMAS record on the role of indigenous languages in Sarawak schools. \url{https://ir.unimas.my/id/eprint/17155/}
\end{itemize}

\section{Linguistic and Reference Materials}
\begin{itemize}
\item Omniglot page on Bukar-Sadong. It states that a Latin-script writing system was devised in 2005. \url{https://omniglot.com/writing/bukar-sadong.htm}
\item Borneo Dictionary, Bidayuh Serian word list. Its quality and size were not checked. \url{https://borneodictionary.com/bidayuh-serian/}
\item Numerals page for Serian Bidayuh at the MPI-EVA Numerals in Languages of the World project. \url{https://lingweb.eva.mpg.de/channumerals/Bidayah-Serian.htm}
\item UNIMAS record on Bidayuh language use and proficiency in Sarawak. \url{https://ir.unimas.my/id/eprint/50111/}
\end{itemize}

\section{Religious and Community Resources}
Nothing specific to Bukar-Sadong scripture or community organisations was confirmed. The Dayak Bidayuh National Association is reported to run Bidayuh Heritage Schools. See the Borneo Post article above and the community-oriented pages below.
\begin{itemize}
\item Bidayuh community information page. Its content was not verified. \url{https://bidayuh.yolasite.com/}
\item Borneo Post, ``Bid to revive Bidayuh language in schools''. \url{https://www.theborneopost.com/?p=77460}
\end{itemize}

\section{Digital and Audio Resources}
\begin{itemize}
\item Wikitongues language page for Bukar-Sadong Bidayuh. It may link to oral recordings; this was not checked. \url{https://wikitongues.org/languages/sdo/}
\end{itemize}
No dedicated audio corpus or software was found.

\section{Curriculum Development Context}
Bidayuh had no standardised orthography until the early 2000s. Reports describe dialect variation as the main barrier to a common school subject. Heritage-language kindergartens and playschools exist, and there is a 2025 plan for an optional school subject. A K-12 programme in Bukar-Sadong would need new materials, a dialect decision and a check of the 2005 orthography with the community.
\end{document}